\section{Conclusions}
\label{sec:conclusions}

We evaluated seven machine learning models for micropollutant adsorption
on layered double hydroxides under a leakage-aware protocol that
separates within-system interpolation from cross-study generalization,
using 1,342 experiments compiled from 59 studies. The models predict
removal ratio for studied systems at new operating conditions with test
$R^2$ of $0.85 \pm 0.03$ and RMSE near 11 percentage points across 20
repeated splits, but the same models show near-zero per-study skill when
entire studies are held out, and a leave-one-study-out audit locates the
failure in a wide per-study error distribution rather than a uniform
degradation. Against descriptor-free baselines the picture is more
nuanced: the models reduce the naive cross-study error by roughly 40
percent, so they do carry transferable signal, yet not enough to certify
predictions for an unmeasured system. Split-conformal
intervals certify interpolation predictions at $\pm$19 percentage points,
whereas for unseen studies neither conformal calibration nor a
distance-based applicability score yields informative guarantees.
Chemistry-informed features such as pH$-$pH$_{\mathrm{pzc}}$ leave the
accuracy essentially unchanged but anchor the learned behavior in
electrostatic and hydrophobic mechanisms. In
future work, we plan to extend the compilation with interlayer-anion and
morphology descriptors, model study-level effects explicitly, and test
few-shot calibration so that a small number of measurements of a new
system can recover the accuracy currently reserved for studied ones.
